\section{Explicit Frey curves realizing the local fields}
\label{sec:frey-realizations}

The local calculations above use actual finite extensions of $\Q_p$.
We now realize their ramification by torsion fields of explicit rational
elliptic curves. This separates a local model from a claim that every
condition on global initial data has already been supplied. The latter
construction is given separately in Section~\ref{sec:initial-theta-data}.

For positive coprime $a,b$, put $c=a+b$ and
\[
 D_{a,b}:\quad y^2=x(x-a)(x+b),\qquad S=a^2+ab+b^2.
\]
Direct calculation from the Weierstrass coefficients gives
\begin{equation}\label{eq:fr-invariants}
 \Delta(D_{a,b})=16(abc)^2,\qquad c_4(D_{a,b})=16S,
 \qquad j(D_{a,b})=\frac{256S^3}{(abc)^2}.
\end{equation}
Reducing $S$ modulo $a,b,c$ gives $b^2,a^2,a^2$, respectively;
thus $\gcd(S,abc)=1$. At an odd prime $r\mid abc$, the displayed
equation is minimal and multiplicative, and its Tate order is
$2v_r(abc)$ after a splitting extension. Indeed $c_4$ is a unit,
which both proves minimality and gives the reduction criterion
\cite[VII.5.1]{Silverman}.

\begin{lemma}[Two transvections and passage to a torsion field]
\label{lem:fr-sl2}
Let $\ell$ be a prime and let $G\subseteq\GL_2(\mathbf F_\ell)$
be a subgroup.
If $G$ contains a nonidentity transvection and an element preserving
no $\mathbf F_\ell$-line, then $G$ contains
$\operatorname{SL}_2(\mathbf F_\ell)$.
If $H\triangleleft G$ has finite index prime to $\ell$, then $H$
also contains this special linear group.
\end{lemma}
\begin{proof}
Let $T$ be the transvection, with fixed line $L$, and let $g$ be the
second element. The lines $L$ and $gL$ differ. In a basis along them,
$T$ and $gTg^{-1}$ have matrices
\[
 U(s)=\begin{pmatrix}1&s\\0&1\end{pmatrix},\qquad
 V(t)=\begin{pmatrix}1&0\\t&1\end{pmatrix},\qquad st\ne0.
\]
Their integer powers give $U(x)$ and $V(y)$ for every
$x,y\in\mathbf F_\ell$, since the field is prime. Elementary
row reduction shows that these two root groups generate
$\operatorname{SL}_2(\mathbf F_\ell)$: for a matrix with nonzero
upper-left entry, a lower row operation first clears the lower-left
entry, and an upper operation clears the upper-right entry; diagonal
matrices are products of elementary matrices. Explicitly,
$W(t)=U(t)V(-t^{-1})U(t)$ has entries
$\left(\begin{smallmatrix}0&t\\-t^{-1}&0\end{smallmatrix}\right)$,
so $W(t)W(-1)=\operatorname{diag}(t,t^{-1})$.
If the upper-left
entry is zero, a preliminary upper operation makes it nonzero.
For the last assertion, the image of $T$ in $G/H$ has order dividing
both $\ell$ and $[G:H]$, hence is trivial. Normality retains
$gTg^{-1}$, so the same argument applies inside $H$.
\end{proof}

Here and below the arithmetic use of this lemma invokes two classical
facts. Good-reduction Frobenius has characteristic polynomial
$T^2-a_rT+r$, with $a_r=r+1-\#D(\mathbf F_r)$
\cite[V.2; VII.4]{Silverman}. At a split multiplicative prime
$p\ne\ell$ with $\ell\nmid v_p(q)$, the Tate description
$\overline{\Q}_p^\times/q^\Z$ supplies a nontrivial inertia
transvection on $D[\ell]$: adjoining $q^{1/\ell}$ has ramification
index $\ell$, while $\mu_\ell$ is unramified. See
\cite[Theorems 1--2 and Proposition 3]{KedlayaTate}.
These statements about arithmetic representations are not consequences
of the formal elementary-matrix lemma alone.

\begin{proposition}[Degree 210 at the prime 139]
\label{prop:fr-139}
Each of the curves
\[
 D_{139,279},\qquad D_{1,2362}
\]
has split multiplicative reduction at $139$ with native Tate order
two. For either curve $D$, the mod-$7$ image over
$F=\Q(i,D[30])$ contains $\operatorname{SL}_2(\mathbf F_7)$.
The completion at $139$ of $\Q(i,D[210])$ is
\begin{equation}\label{eq:fr-210-field}
 E=\Q_{139}(\mu_{210},\pi),\qquad
 \pi^{105}=b_0,\qquad b_0^2=q,\quad v_{139}(b_0)=1.
\end{equation}
It has ramification index $105$, residue degree two, and degree $210$.
The second curve is already a rational Legendre curve, with
parameter $-2362$.
\end{proposition}
\begin{proof}
The factorizations are
\[
 279=3^2\cdot31,\quad 418=2\cdot11\cdot19,
 \qquad2362=2\cdot1181,\quad2363=17\cdot139.
\]
All displayed factors other than $3^2$ are prime.
For the first curve, reduction at $139$ gives $y^2=x^2(x+1)$;
for the second it gives $y^2=x(x-1)^2$. The respective nodes have
two distinct rational tangent directions, of slopes $\pm1$.
Equation~\eqref{eq:fr-invariants} gives $v_{139}(\Delta)=2$
and unit $c_4$ in both cases.

Counting the actual points modulo five, including infinity, gives
$8$ and $4$, respectively. The two Frobenius polynomials are
$T^2+2T+5$ and $T^2-2T+5$; both have discriminant $-16$, a
nonsquare modulo seven. The transvection at $139$ and
Lemma~\ref{lem:fr-sl2} therefore give the special linear subgroup
over $\Q$. The field $F$ is Galois and
\begin{equation}\label{eq:fr-level-degree}
 [F:\Q]\mid2\,|\GL_2(\Z/30\Z)|
       =2\cdot6\cdot48\cdot480=276480=2^{11}3^3\cdot5.
\end{equation}
Its degree is prime to seven, so the normal-subgroup part of the
lemma applies over $F$.

Full rational two-torsion makes a square root $b_0$ of $q$ rational
over $\Q_{139}$. In fact its Tate class is rational, so for every
Galois element $\sigma$ the quotient $\sigma(b_0)/b_0$ belongs
to both $\{\pm1\}$ and $q^\Z$. Its valuation is zero, forcing
that quotient to be one. The order of $139$ modulo $210$ is two,
so the cyclotomic field in \eqref{eq:fr-210-field} is unramified
quadratic. The polynomial $T^{105}-b_0$ is Eisenstein there.
The Tate torsion classes give
$\Q_{139}(D[210])=\Q_{139}(\mu_{210},q^{1/210})=E$;
fixing their classes forces the root-of-unity multipliers to be
one. The unramified quadratic field also contains $i$. Finally
$D_{1,2362}$ is $y^2=x(x-1)(x+2362)$, as asserted.
No equality $q=139^2$ has been assumed.
\end{proof}

These small examples require care if used with the particular
finite-exception construction of \cite[\S5.8]{JoshiIV}. Write
\[
 N_1=8105229=139\cdot279\cdot418/2,
 \qquad N_2=2790703=2362\cdot2363/2.
\]
In both cases the normalized Tate quantity is $Q_D=2\log N_i$.
At two, $v_2(j)=6$; adjoining the full three-torsion gives good
reduction there, by the finite-torsion argument below. Summing the
odd pole orders of $j$ gives the asserted normalized value.
The exact inequalities
\[
 3^{25}<N_i^2<2^{49},\qquad i=1,2,
\]
and $2<\exp(1)<3$ imply $25<Q_D<49$. But the prime threshold
in that existence proof must exceed ten, since
$\theta(10)=\log210<6<20/3$. Thus candidates with
$\sqrt{Q_D}\le7$ belong to its proof-constructed small-height
exceptional set when they lie in the chosen bounding domain.
This observation uses the normalized definition of $Q_D$ in
Definition~5.4.1 and Theorem~5.7.1 of that source. The opening
display of its Section~5.8 omits the degree divisor; a literal
unnormalized reading must not reuse the preceding normalized values.
This is a limitation of invoking that sufficient existence proof,
not a failure of the local calculations or of a direct initial-data
construction.

\begin{theorem}[An exact balanced example at level 43]
\label{thm:fr-balanced}
Put
\[
 p=1289,\quad\ell=43,\quad A=p(p^{16}+428),\qquad
 (a,b,c)=(A^2,A^2+1,2A^2+1),\quad D=D_{a,b}.
\]
This is a primitive abc triple, and $D$ is rationally isomorphic
to the Legendre curve of parameter $\lambda=-1-A^{-2}$.
For $F=\Q(i,D[30])$, the following assertions hold:
\begin{enumerate}
\item $D$ is good above two and at $43$, and split multiplicative
at every odd prime dividing $abc$.
\item The mod-$43$ image over both $\Q$ and $F$ contains
$\operatorname{SL}_2(\mathbf F_{43})$.
\item Every nonzero integral Tate order over $F$ is prime to $43$.
\item The normalized Tate quantity satisfies the exact bounds
\begin{equation}\label{eq:fr-balanced-height}
 Q_D=2\log\left(\frac{abc}{2}\right),
       \qquad 1224<Q_D<1648<43^2.
\end{equation}
\item At $p$, the completion of $K=\Q(i,D[1290])$ has degree
$1290$, ramification index $645$, and residue degree two.
\end{enumerate}
\end{theorem}
\begin{proof}
The three endpoints are pairwise coprime, and $a+b=c$.
The rational substitution $x=A^2X$, $y=A^3Y$ gives
\[
 Y^2=X(X-1)(X+1+A^{-2}),
\]
so no quadratic twist is introduced. The exact residues are
$A\equiv5\pmod8$ and $A\equiv1\pmod{5\cdot43}$.
At an odd prime dividing $a$ or $b$, the reduced equation is
$y^2=x^2(x+1)$, with tangent slopes $\pm1$.
At a prime dividing $c$, translation of its node gives
$y^2=X^2(X+A^2)$, with slopes $\pm A$.
All these slopes are nonzero. The endpoints are $(1,2,3)$
modulo $43$, so the discriminant is a unit there.

At two, $v_2(abc)=1$ and $3A^4+3A^2+1$ is odd, whence
$v_2(j)=6$. Here is a direct reason that the particular field
$F$ gives good reduction. At any completion $F_w$ above two,
the group $D[3]\simeq(\Z/3\Z)^2$ is rational.
If reduction were additive, both the formal subgroup $D_1(F_w)$
and the quotient $D_0(F_w)/D_1(F_w)$ would have no three-torsion.
Thus $D[3]$ would inject into $D(F_w)/D_0(F_w)$, whose order
is at most four in the additive case, a contradiction.
These are \cite[VII.2.1, VII.3.1, VII.5.1(c), VII.6.1]{Silverman},
valid over finite extensions of $\Q_2$ as well.
Multiplicative reduction would require negative $j$-valuation.
The only remaining possibility is good reduction.

At five the reduced model has four points and Frobenius trace two.
The discriminant $-16$ of $T^2-2T+5$ is a nonsquare modulo $43$.
Also $v_p(A)=1$, giving native Tate order four at $p$.
Lemma~\ref{lem:fr-sl2}, followed by
\eqref{eq:fr-level-degree}, proves the image assertion.

For clarity, the assertion about all prime exponents does not use
a factorization of the large endpoints. Let
$P=\prod_{r\le511,\ r\text{ prime}}r$ be the actual primorial.
The exact integer calculations give
\[
 a,b,c<2^{377},\qquad
 \gcd(a,P)=3,\quad\gcd(b,P)=2,\quad\gcd(c,P)=17,
\]
and $v_3(a)=2$, $v_2(b)=1$, $v_{17}(c)=1$.
These identities can be verified by integer Euclidean division;
the companion Lean proof also checks the primorial defined as the
product over all primes in the indicated range.
Every prime $r\ge512$ has $r^{43}\ge2^{387}>a,b,c$.
Since the endpoints are pairwise coprime, it follows that
\begin{equation}\label{eq:fr-all-primes}
 v_r(abc)<43\qquad\hbox{for every prime }r.
\end{equation}
At a place $w\mid r$ of $F$, the nonzero integral Tate order is
$2e(w/r)v_r(abc)$. The ramification index divides the Galois
degree in \eqref{eq:fr-level-degree}, which is prime to $43$;
equation~\eqref{eq:fr-all-primes} proves the claim.

The normalized sum over $F$ is the sum of the rational odd
pole orders, because
\[
 \sum_{w\mid r}e(w/r)f(w/r)=[F:\Q].
\]
This proves the identity in \eqref{eq:fr-balanced-height}.
For its strict bounds, put $N=abc/2$. We have
$p^{17}<A<2p^{17}$ and $A^6<N<3A^6$.
The integer inequalities $3^6<p<(8/3)^8$ give
\[
 3^{1224}<N^2
       <9\cdot2^{12}(8/3)^{1632}<(8/3)^{1648}.
\]
For the last step use
$9\cdot2^{12}<3^{12}<(8/3)^{16}$.
Together with $8/3<\exp(1)<3$, these imply
\eqref{eq:fr-balanced-height} without a floating point approximation.

Finally full rational two-torsion gives $b_0=\sqrt q\in\Q_p$
with $v_p(b_0)=2$, by the proof of Proposition~\ref{prop:fr-139}.
Put $K_0=\Q_p(\mu_{1290})$ and choose $\pi^{645}=b_0$.
The field $K_0$ is unramified quadratic, since $p\equiv-1\pmod{1290}$.
Writing $\eta=p^2/b_0\in\Z_p^\times$, the genuine uniformizer is
\begin{equation}\label{eq:fr-uniformizer}
 \beta=\pi^{323}/p,\qquad
 \beta^{645}=p\eta^{-323},\qquad\pi=\eta\beta^2.
\end{equation}
The first polynomial in $\beta$ is Eisenstein. Tate torsion
identifies $K_v=K_0(\pi)=K_0(\beta)$, and $K_0$ contains $i$.
The claimed degrees follow. In particular the native root
$q^{1/86}=\pi^{15}$ has valuation $2/43$, not $1/43$.
\end{proof}

The numerical upper bound on the level in
\cite[Theorem~5.7.1]{JoshiIV} is also satisfied here: with
$\delta=2^{12}3^3\cdot5$ and $d_{\rm mod}=1$,
\[
 \sqrt{Q_D}<43
       <10\delta\sqrt{Q_D}\log(2\delta\log Q_D).
\]
This directly verifies a window for the specified prime. It does
not identify that prime with an unspecified existential choice in
another theorem. The next sections construct the additional data
and an unbounded family independently of that identification.
